\documentclass[a4paper,12pt]{article}
\usepackage[utf8]{inputenc}
\usepackage[T1]{fontenc}
\usepackage[spanish]{babel}
\usepackage{enumitem}
\usepackage{geometry}
\usepackage{hyperref}
\usepackage{titlesec}
\usepackage{array}

\geometry{margin=2.5cm}

\titleformat{\section}{\large\bfseries}{\thesection}{1em}{}
\titleformat{\subsection}{\normalsize\bfseries}{\thesubsection}{1em}{}

\title{Flujo completo de componentes --- Madness Dreams}
\author{Donovan Eliam Aguilar Diaz}
\date{\today}

\begin{document}

\maketitle
\tableofcontents
\newpage

\section{Visión general}

El proyecto se organiza en capas:

\begin{enumerate}
    \item \textbf{Datos (ScriptableObjects)}: configuración editable sin lógica de frame.
    \item \textbf{Generación procedural}: construye el mapa con semilla.
    \item \textbf{Enemigos}: FSM, movimiento, ataque y visión.
    \item \textbf{Jefes}: FSM con fases y arena.
    \item \textbf{Núcleo}: vida, proyectiles, flujo de mundos y enumeraciones.
    \item \textbf{Jugador}: movimiento, combate y animación.
\end{enumerate}

\subsection{Flujo general}

\texttt{RunManager} carga \texttt{WorldConfigSO}. Luego \texttt{MapGenerator} estampa los chunks con \texttt{TilemapStamper}. Después \texttt{ChunkContent} y \texttt{Spawner} colocan enemigos. El jugador explora y combate. Al entrar en la arena, \texttt{BossArenaController} spawnea al jefe. Al morir, \texttt{RunManager} pasa al siguiente mundo.

\section{Capa de datos (ScriptableObjects)}

\subsection{BlockBankSO}

\begin{itemize}
    \item \textbf{Propósito}: banco de chunks disponibles por nivel y la arena del jefe.
    \item \textbf{Campos}: \texttt{normalChunks} (lista de \texttt{MapChunkSO}), \texttt{bossArena}.
    \item \textbf{Métodos}: \texttt{GetRandomChunk(System.Random)}, \texttt{GetChunk(int)}.
    \item \textbf{Flujo}: \texttt{MapGenerator} pide un chunk aleatorio por celda del grid.
\end{itemize}

\subsection{MapChunkSO}

\begin{itemize}
    \item \textbf{Propósito}: plantilla de un bloque o room (tiles, puertas y puntos de spawn).
    \item \textbf{Campos}: \texttt{size}, \texttt{tiles} (aplanado fila a fila), \texttt{doorPositions}, \texttt{enemySpawns}, \texttt{role}.
    \item \textbf{Helpers}: \texttt{GetTile(x, y)}, \texttt{LocalArea}, \texttt{SetSize}, \texttt{SetTileData}.
    \item \textbf{Flujo}: el editor captura un \texttt{Tilemap} pintado con \texttt{GetTilesBlock} y \texttt{TilemapStamper} lo estampa con \texttt{SetTilesBlock}.
\end{itemize}

\subsection{EnemyStatsSO}

\begin{itemize}
    \item \textbf{Propósito}: plantilla de propiedades comunes de un enemigo.
    \item \textbf{Campos}: \texttt{moveSpeed}, \texttt{chaseSpeed}, \texttt{maxHealth}, \texttt{damage}, \texttt{attackRange}, \texttt{attackCooldown}, \texttt{windupTime}, \texttt{attackType}, \texttt{projectilePrefab}, \texttt{projectileSpeed}, \texttt{detectionRadius}, \texttt{fovRadius}, \texttt{fovAngle}, \texttt{hearingRadius}.
    \item \textbf{Flujo}: \texttt{EnemyAI} inyecta estos stats en \texttt{EnemyMovement}, \texttt{EnemyAttack}, \texttt{EnemyFOV} y \texttt{Health}.
\end{itemize}

\subsection{BossStatsSO}

\begin{itemize}
    \item \textbf{Hereda de}: \texttt{EnemyStatsSO}.
    \item \textbf{Campos extra}: \texttt{phaseHealthThresholds}, \texttt{phaseAttacks}.
    \item \textbf{Flujo}: \texttt{BossAI} usa los umbrales para detectar cambios de fase.
\end{itemize}

\subsection{EnemyArchetypeSO}

\begin{itemize}
    \item \textbf{Propósito}: vincula un prefab con sus stats.
    \item \textbf{Campos}: \texttt{prefab}, \texttt{stats}.
\end{itemize}

\subsection{WorldConfigSO}

\begin{itemize}
    \item \textbf{Propósito}: configuración de un mundo.
    \item \textbf{Campos}: \texttt{seed}, \texttt{bank}, \texttt{width}, \texttt{height}, \texttt{chunkSize}, \texttt{bossRegionSize}, \texttt{enemyPool}, \texttt{enemiesPerChunk}, \texttt{bossPrefab}, \texttt{boss}.
\end{itemize}

\subsection{MapChunkCaptureEditor (editor)}

\begin{itemize}
    \item \textbf{Propósito}: capturar un room pintado a un \texttt{MapChunkSO}.
    \item \textbf{Flujo}: \texttt{Tools > Madness Dreams > Map Chunk Capture}, asignar Tilemap y target, clic en \texttt{Capture}.
\end{itemize}

\section{Generación procedural}

\subsection{CellRole}

Enumeración: \texttt{Empty}, \texttt{Start}, \texttt{Normal}, \texttt{Boss}. Marca el rol de cada celda del grid.

\subsection{TilemapStamper}

\begin{itemize}
    \item \textbf{Propósito}: estampar un \texttt{MapChunkSO} en un Tilemap global.
    \item \textbf{Flujo}: \texttt{Stamp(tilemap, chunk, origin, rotationSteps)} rota el layout si es necesario y llama a \texttt{SetTilesBlock}.
\end{itemize}

\subsection{MapGenerator}

\textbf{Propósito}: orquestar toda la generación.

\textbf{Flujo}:

\begin{enumerate}
    \item Valida \texttt{config}, \texttt{bank} y \texttt{targetTilemap}; limpia el tilemap.
    \item Crea \texttt{System.Random(seed)}.
    \item \texttt{BuildRoles}: \texttt{Start} al centro, \texttt{Boss} en la celda más lejana (región \texttt{bossRegionSize}), resto \texttt{Normal}.
    \item \texttt{StampNormalChunks}: por celda \texttt{Normal} o \texttt{Start}, pide un chunk al banco, rota y estampa.
    \item \texttt{StampBossArena}: estampa la arena en la región del jefe.
    \item \texttt{SpawnEnemies}: delega en \texttt{ChunkContent}.
\end{enumerate}

\subsection{Spawner}

\begin{itemize}
    \item \textbf{Propósito}: instanciar con chequeo de colisión.
    \item \textbf{Flujo}: \texttt{TrySpawn(prefab, position, rng)} genera candidatos con jitter, usa \texttt{Physics2D.OverlapCircle} e instancia si no hay solapamiento (hasta \texttt{maxTries}).
\end{itemize}

\subsection{ChunkContent}

\begin{itemize}
    \item \textbf{Propósito}: poblar un chunk con enemigos.
    \item \textbf{Flujo}: lee \texttt{enemySpawns}, elige enemigos del \texttt{enemyPool}, convierte celda a mundo con \texttt{GetCellCenterWorld} y usa el \texttt{Spawner}.
\end{itemize}

\section{Enemigos}

\subsection{EnemyAI (FSM central)}

\textbf{Estados}: \texttt{Patrol}, \texttt{Alert}, \texttt{Chase}, \texttt{Attack}, \texttt{Return}, \texttt{Dead}.

\textbf{Por frame}: \texttt{attack.Tick}, \texttt{UpdatePerception} (FOV con intervalo), switch de estado y \texttt{UpdateAnimation}.

\textbf{Transiciones}:

\begin{center}
\begin{tabular}{|l|l|}
\hline
\textbf{De a} & \textbf{Condición} \\
\hline
\texttt{Patrol} $\to$ \texttt{Alert} & Jugador visto u oído \\
\hline
\texttt{Alert} $\to$ \texttt{Chase} & Terminó el tiempo de alerta \\
\hline
\texttt{Chase} $\to$ \texttt{Attack} & Distancia $\leq$ \texttt{attackRange} \\
\hline
\texttt{Attack} $\to$ \texttt{Chase} & Jugador salió de \texttt{attackRange} \\
\hline
\texttt{Chase} $\to$ \texttt{Return} & Jugador fuera de \texttt{detectionRadius} por un tiempo \\
\hline
\texttt{Return} $\to$ \texttt{Patrol} & Llegó al punto de origen \\
\hline
Cualquiera $\to$ \texttt{Dead} & Vida $\leq$ 0 \\
\hline
\end{tabular}
\end{center}

\subsection{EnemyMovement}

\texttt{MoveDirection}, \texttt{MoveTowards} y \texttt{Stop} setean la velocidad del \texttt{Rigidbody2D}.

\subsection{EnemyAttack}

Según \texttt{stats.attackType}:

\begin{itemize}
    \item \texttt{Melee}: corrutina con \texttt{windupTime} y luego aplica daño en rango.
    \item \texttt{Ranged}: corrutina con \texttt{windupTime} y luego instancia \texttt{stats.projectilePrefab} y lo lanza.
\end{itemize}

\subsection{EnemyFOV}

\texttt{CanHearPlayer} (radio) y \texttt{CanSeePlayer} (radio, cono y línea de visión con raycast contra \texttt{obstacleMask}).

\subsection{Projectile}

\texttt{Launch} setea dirección y velocidad, vuela con \texttt{Rigidbody2D}, y en \texttt{OnTriggerEnter2D} daña si el collider está en \texttt{targetMask}; se autodestruye por \texttt{lifetime}.

\section{Jefes}

\subsection{BossAI}

\begin{itemize}
    \item \textbf{Hereda de}: \texttt{EnemyAI}.
    \item \textbf{Flujo}: calcula la fase según la vida restante (\texttt{phaseHealthThresholds}) y, en fase 2 o superior, invoca esbirros con el \texttt{Spawner}.
\end{itemize}

\subsection{BossSpawnPoint}

Marcador de posición del jefe (con gizmo).

\subsection{BossArenaController}

\texttt{OnTriggerEnter2D} (capa \texttt{BossTrigger}) bloquea las puertas, spawnea al jefe y, al morir, desbloquea y dispara \texttt{BossDefeated}.

\subsection{BossManager}

Configura la arena desde \texttt{WorldConfigSO} (\texttt{Setup}) y escucha \texttt{BossDefeated} para avisar a \texttt{RunManager}.

\section{Núcleo}

\subsection{Health / IDamageable}

\texttt{TakeDamage} reduce la vida, dispara \texttt{HealthChanged} y, al llegar a 0, \texttt{Died}.

\subsection{RunManager}

En \texttt{Start} suscribe \texttt{BossDefeated} y genera el mundo 0; al derrotar al jefe, avanza al siguiente mundo.

\subsection{AnimationState}

Enumeración: \texttt{Idle}, \texttt{Move}, \texttt{MeleeAttack}, \texttt{RangedAttack}, \texttt{Death}.

\subsection{AttackType}

Enumeración: \texttt{Melee}, \texttt{Ranged}.

\section{Jugador}

\subsection{PlayerController}

Lee el input en \texttt{Update} y setea la velocidad en \texttt{FixedUpdate}; expone \texttt{IsMoving}.

\subsection{PlayerCombat}

\begin{itemize}
    \item \texttt{Melee} (tecla Space): \texttt{OverlapCircleAll} contra \texttt{enemyMask} y daño a cada \texttt{IDamageable}.
    \item \texttt{Ranged} (tecla K): instancia \texttt{projectilePrefab} hacia el ratón y lo lanza.
    \item Expone \texttt{IsMeleeAttacking} e \texttt{IsRangedAttacking}.
\end{itemize}

\subsection{PlayerAnimationController}

Cada frame determina \texttt{AnimationState} (muerte, ataque, movimiento o idle) y setea el parámetro \texttt{State} del Animator.

\section{Flujo end-to-end}

\begin{enumerate}
    \item \texttt{RunManager} carga el mundo 0.
    \item \texttt{MapGenerator} genera tiles y reserva la región del jefe.
    \item \texttt{ChunkContent} y \texttt{Spawner} spawnean esbirros.
    \item El jugador explora; los enemigos detectan, persiguen y atacan.
    \item Al entrar en la arena, \texttt{BossArenaController} sella las puertas y spawnea al jefe.
    \item \texttt{BossAI} pelea cambiando de fase.
    \item El jefe muere y \texttt{BossManager}/\texttt{RunManager} avanzan al siguiente mundo con nueva semilla.
\end{enumerate}

\section{Resumen de estados de animación}

Parámetro entero \texttt{State} del Animator:

\begin{center}
\begin{tabular}{|c|l|}
\hline
\textbf{Valor} & \textbf{Estado} \\
\hline
0 & \texttt{Idle} \\
1 & \texttt{Move} \\
2 & \texttt{MeleeAttack} \\
3 & \texttt{RangedAttack} \\
4 & \texttt{Death} \\
\hline
\end{tabular}
\end{center}

\subsection{Mapeo enemigo (FSM a animación)}

\begin{center}
\begin{tabular}{|l|l|}
\hline
\textbf{Estado FSM} & \textbf{Animación} \\
\hline
\texttt{Patrol} (quieto) & \texttt{Idle} \\
\texttt{Patrol} (moviendo) & \texttt{Move} \\
\texttt{Alert} & \texttt{Idle} \\
\texttt{Chase} & \texttt{Move} \\
\texttt{Return} & \texttt{Move} \\
\texttt{Attack} (melee) & \texttt{MeleeAttack} \\
\texttt{Attack} (ranged) & \texttt{RangedAttack} \\
\texttt{Dead} & \texttt{Death} \\
\hline
\end{tabular}
\end{center}

\subsection{Mapeo jugador}

\begin{center}
\begin{tabular}{|l|l|}
\hline
\textbf{Condición} & \textbf{Animación} \\
\hline
Vida $\leq$ 0 & \texttt{Death} \\
\texttt{IsMeleeAttacking} & \texttt{MeleeAttack} \\
\texttt{IsRangedAttacking} & \texttt{RangedAttack} \\
\texttt{IsMoving} & \texttt{Move} \\
resto & \texttt{Idle} \\
\hline
\end{tabular}
\end{center}

\end{document}
